Retrieval-Augmented Generation (RAG) has emerged as a dominant paradigm for grounding large language model outputs in factual, domain-specific knowledge. 
Yet practitioners face a vast, poorly understood configuration space spanning chunking strategies, embedding models, retrieval methods, and generation models. 
We present a modular, extensible benchmarking framework that systematically evaluates RAG pipelines across this configuration space. 
Our framework orchestrates the full RAG lifecycle - document chunking, embedding, retrieval strategies, generation, and evaluation — supporting distributed execution across heterogeneous GPU infrastructure. 
Evaluation combines RAGAS quality metrics with information retrieval metrics and natural language generation metrics. 
The framework is demonstrated on Wikipedia QA datasets, benchmarking over 100 configuration combinations across locally-served open-weight models. 
The results reveal non-obvious interactions between chunking granularity, retrieval strategy, and generation quality. 
The framework supports extensible dataset adapters, per-role endpoint configuration for distributed setups, and integrated experiment tracking via Mlflow.